% =============================================================================
% 036 / 068
% Status: raw
% =============================================================================
% Notes:
%
% Source question:
% 目前，很少見到你著墨於構詞法(造詞的原則)？
% 閣下的文章幾乎都是以梳理語法為主。
% 
% 雖然語法重要性很大，
% 但是構詞法、辨別原生吳語詞、保留非藍青的詞彙或句型，
% 重要程度不遜於語法呢?
% =============================================================================

\chapter[目前，很少見到你著墨於構詞法(造詞的原則)？ 閣下的文章幾乎都是以梳理語法為主。 雖然語法重要性很大， 但是構詞法、辨別原生吳語詞、保留非藍青的詞彙或句型， 重要程度不遜於語法呢?]{目前，很少見到你著墨於構詞法(造詞的原則)？ \\[0.35em] 閣下的文章幾乎都是以梳理語法為主。 \\[0.35em] 雖然語法重要性很大， \\[0.35em] 但是構詞法、辨別原生吳語詞、保留非藍青的詞彙或句型， \\[0.35em] 重要程度不遜於語法呢?}
\qaanswer{
詞彙需要編纂『漢吳辭典』，這件事即使有結果在當下的貴國也是無法公開傳播的。
}

